%=====================================================================
% PROJETO INDUSTRIAL II — PLANTA COM SUBESTACAO PROPRIA E SPDA
% Subestação externa cercada, sala elétrica adjacente e fluxo logístico.
%=====================================================================
\section{Projeto Industrial II — empreendimento com subestação própria}

\newcommand{\IndustriaSEFinalBase}{%
\draw[black!25,dashed] (0,0) rectangle (70,45);
\draw[arqExt] (10,0) rectangle (70,45);
% faixa técnica superior
\ArqParedeInt{10}{34}{70}{34}
\foreach \x in {16,28,38,48,58}{\ArqParedeInt{\x}{34}{\x}{45}}
% processo e armazenagem
\foreach \x in {22,36,50,60}{\ArqParedeInt{\x}{7}{\x}{34}}
\ArqParedeInt{10}{7}{70}{7}
% subestação externa e gerador: áreas restritas com acesso próprio
\draw[corAt!75,dashed,line width=0.9pt] (0.6,30.2) rectangle (9.2,44.4);
\ArqPortaCorrerH{3.2}{5.6}{30.2}
\draw[arqEquip,line width=0.8pt] (1.4,31.2) rectangle (8.4,43.4);
\ArqPortaVRight{8.4}{35.2}{0.90}
\ArqMaquina{2.0}{38.5}{5.8}{3.4}{CÉLULAS MT}\ArqMaquina{2.0}{33.0}{5.8}{4.2}{TRAFO 500 kVA}
\node[arqLabel,text=corAt!80!black] at (4.9,43.9){S.E.};
\draw[corAt!50,dashed] (0.8,20.5) rectangle (8.8,28.2);
\ArqPortaCorrerH{3.2}{5.6}{20.5}\ArqMaquina{1.5}{22.0}{6.6}{4.5}{GERADOR}
% faixa técnica: cada ambiente possui acesso próprio pelo corredor interno
\draw[arqEquip] (10.5,34.5) rectangle (15.5,44.5);\node[arqLabel] at (13,43.7){ELÉTRICA};\node[arqSub] at (13,39.5){QGBT / CCM};
\ArqMaquina{17.0}{36.0}{4.2}{2.4}{AR COMP.}\ArqMaquina{22.0}{36.0}{4.2}{2.4}{BOMBAS}\node[arqLabel] at (22,43.7){UTILIDADES};
\ArqMaquina{29.5}{36.0}{3.2}{2.4}{CHILLER}\ArqMaquina{34.0}{36.0}{2.5}{2.4}{TORRE}\node[arqLabel] at (33,43.7){FRIO};
\ArqMaquina{39.5}{36.0}{3.0}{1.8}{OFICINA}\ArqMaquina{43.0}{36.0}{3.0}{1.8}{PEÇAS}\node[arqLabel] at (43,43.7){MANUT.};
\ArqDesk{49.0}{36.0}{1.6}{0.6}\ArqDesk{52.0}{39.0}{1.6}{0.6}\node[arqLabel] at (53,43.7){ADMIN./PCP};
\ArqArmario{59.0}{35.0}{1.0}{2.4}\ArqMesa{62.0}{37.0}{4.0}{1.2}\node[arqLabel] at (64,43.7){VEST./REF.};
\ArqPortaHUp{11.5}{34}{0.90}\ArqPortaHUp{18.0}{34}{0.90}\ArqPortaHUp{30.0}{34}{0.90}\ArqPortaHUp{40.0}{34}{0.90}\ArqPortaHUp{50.0}{34}{0.90}\ArqPortaHUp{60.0}{34}{0.90}
\ArqPortaHDown{11.7}{45}{0.90}
% acessos de pedestres independentes das docas
\ArqPortaHUp{23.0}{0}{1.00}\ArqPortaHUp{66.0}{0}{1.00}
\node[arqSub] at (23.5,1.0){ACESSO PEDESTRES};\node[arqSub] at (66.5,1.0){SAÍDA DE SERVIÇO};
% acessos do corredor logístico aos setores
\ArqVaoH{12.5}{18.0}{7}\ArqVaoH{26.0}{31.0}{7}\ArqVaoH{40.0}{45.0}{7}\ArqVaoH{53.0}{57.0}{7}\ArqVaoH{63.0}{67.0}{7}
% passagens entre etapas para material/equipamento
\foreach \x in {22,36,50,60}{\ArqVaoV{15.5}{19.5}{\x}}
% recebimento/MP
\foreach \y in {9,14,19,24,29}{\draw[arqMob] (10.8,\y) rectangle (21.2,\y+1.5);}\node[arqLabel] at (16,33.2){RECEBIMENTO / MP};
% processo 1
\ArqMaquina{23.5}{10}{4.0}{2.3}{P1-A}\ArqMaquina{30.0}{10}{4.0}{2.3}{P1-B}\ArqMaquina{23.5}{21}{4.0}{2.3}{P1-C}\ArqMaquina{30.0}{21}{4.0}{2.3}{P1-D}\node[arqLabel] at (29,33.2){PROCESSO 1};
% processo 2
\ArqMaquina{37.5}{10}{4.0}{2.3}{P2-A}\ArqMaquina{44.0}{10}{4.0}{2.3}{P2-B}\ArqMaquina{37.5}{21}{4.0}{2.3}{P2-C}\ArqMaquina{44.0}{21}{4.0}{2.3}{P2-D}\node[arqLabel] at (43,33.2){PROCESSO 2};
% embalagem e PA
\foreach \y in {11,17,23,29}{\ArqMaquina{52.0}{\y}{6.0}{1.5}{LINHA EMB.}}\node[arqLabel] at (55,33.2){EMBALAGEM};
\foreach \y in {9,14,19,24,29}{\draw[arqMob] (61.0,\y) rectangle (69.2,\y+1.5);}\node[arqLabel] at (65,33.2){ESTOQUE PA};
% docas: grandes vãos reservados ao fluxo de carga
\ArqPortaCorrerH{12}{20}{0}\node[arqSub] at (16,1.0){DOCA RECEB.};
\ArqPortaCorrerV{9}{13}{70}\ArqPortaCorrerV{17}{21}{70}\ArqPortaCorrerV{25}{29}{70}\node[arqSub,rotate=90] at (68.8,18){DOCAS EXP.};
% fluxo
\draw[-{Stealth[length=3mm]},very thick,corNorma] (15,17)--(29,17)--(43,17)--(55,17)--(65,17)--(69.5,17);
\node[arqSub] at (40,5.5){CORREDOR LOGÍSTICO};
% MT subterrânea para sala elétrica
\draw[corFix!80,dashed,very thick] (8.4,36)--(10.5,39.5);\node[arqSub,text=corFix] at (8.8,38.5){MT};
\CotaH{10}{70}{45.22}{60,00 m nave}\CotaV{0}{45}{70.22}{45,00 m lote}
}

\begin{frame}{Industrial II — acessos técnicos, logísticos e de pessoas}
\centering
\begin{tikzpicture}[scale=0.145,every node/.style={font=\tiny}]\IndustriaSEFinalBase\end{tikzpicture}
\vspace{0.4mm}\scriptsize Subestação e gerador possuem acessos externos restritos; a faixa técnica tem portas próprias; pedestres entram por portas dedicadas e as docas permanecem exclusivas ao fluxo logístico.
\end{frame}

\begin{frame}{Industrial II — distribuição física de força}
\centering
\begin{tikzpicture}[scale=0.145,every node/.style={font=\tiny}]
\IndustriaSEFinalBase
\SimQD{13}{39.5}{QGBT}\SimQD{15}{37.0}{CCM}
\foreach \x/\y/\n in {25.5/11/M1,32/11/M2,25.5/22/M3,32/22/M4,39.5/11/M5,46/11/M6,39.5/22/M7,46/22/M8}{\SimMotor{\x}{\y}{\n}}
\foreach \x/\y in {53/13,57/13,53/25,57/25}{\SimTomadaTri{\x}{\y}}
\foreach \x in {14,20,26,32,38,44,50,56,62,68}{\SimLuz{\x}{9.5}\SimLuz{\x}{28}}
\RotaF{15}{37}{25.5}{22}\RotaF{15}{37}{39.5}{22}\RotaF{15}{37}{55}{25}
\draw[corPrim!70,dashed,thick] (13,39.5)--(20,32)--(60,32);
\end{tikzpicture}
\end{frame}

\begin{frame}{Industrial II — demanda e transformador}
\small
\begin{columns}[T,onlytextwidth]
\column{0.49\textwidth}\begin{caixaProjeto}[title=Inventário do estudo]\begin{itemize}\item processo 1: 120 kW;\item processo 2: 105 kW;\item embalagem: 42 kW;\item refrigeração/utilidades: 72 kW;\item iluminação/TUG/admin.: 35 kW;\item reservas e auxiliares: 45 kW.\end{itemize}\end{caixaProjeto}
\column{0.49\textwidth}\begin{caixaNorma}[title=Resultado preliminar]\begin{itemize}\item instalada: cerca de 419 kW;\item demanda: cerca de 340 kW;\item $fp$ de estudo: 0,92;\item $S_d\approx370$ kVA;\item trafo didático: 500 kVA.\end{itemize}\end{caixaNorma}
\end{columns}
\vspace{1mm}\begin{caixaAt}A escolha final do transformador e da entrada em MT depende da demanda, expansão, partida de motores, qualidade de energia e análise técnica da Equatorial.\end{caixaAt}
\end{frame}

\begin{frame}{Industrial II — subestação e unifilar MT/BT}
\centering
\begin{tikzpicture}[scale=0.69,every node/.style={font=\tiny}]
\tikzset{b/.style={draw=corPrim,fill=corDef!8,rounded corners,minimum width=2.2cm,minimum height=0.58cm,align=center}}
\node[b](r)at(0,5.0){Rede MT Equatorial};\node[b](med)at(0,4.1){Medição / proteção MT};\node[b](tr)at(0,3.2){Trafo 500 kVA\\13,8 kV / 380-220 V*};\node[b](qg)at(0,2.3){QGBT 1.000 A};
\node[b](ccm)at(-4.2,0.9){CCM processo};\node[b](uti)at(-1.4,0.9){QD utilidades};\node[b](emb)at(1.4,0.9){QD embalagem};\node[b](ser)at(4.2,0.9){QD serviços};
\draw[-{Stealth},thick,corPrim](r)--(med)--(tr)--(qg);\foreach \x in {ccm,uti,emb,ser}{\draw[-{Stealth},thick](qg)--(\x);}
\draw[corNorma,thick](qg.east)--(5.8,2.3)--(5.8,-0.3)node[ground]{};\node[corNorma,anchor=west] at (5.85,1.2){BEP / malha geral};
\node[font=\tiny,anchor=west] at (-5.6,-0.2){* confirmar tensão primária no ponto de conexão real.};
\end{tikzpicture}
\end{frame}

\begin{frame}{Industrial II — curto-circuito preliminar}
\small
\[I_{N,sec}=\frac{500000}{\sqrt3\cdot380}\approx760\,\mathrm{A}\]
Com impedância didática $Z=5{,}5\%$:
\[I_{cc,trafo}\approx\frac{760}{0{,}055}\approx13{,}8\,\mathrm{kA}\]
\begin{caixaAt}Este não é o curto final do QGBT. O estudo executivo incorpora rede MT, transformador real, cabos/barramentos, conexões e contribuição de motores antes de definir Icu/Ics e seletividade.\end{caixaAt}
\end{frame}

\begin{frame}{Industrial II — SPDA e aterramento integrados}
\centering
\begin{tikzpicture}[scale=0.135,every node/.style={font=\tiny}]
\draw[arqExt] (10,0) rectangle (70,45);
\draw[corAt!70,dashed] (0.6,30.2) rectangle (9.2,44.4);\node[arqLabel,text=corAt!80!black] at (4.9,42.8){S.E.};
\foreach \y in {0,9,18,27,36,45}{\draw[corAt!78,thick] (10,\y)--(70,\y);}
\foreach \x in {10,20,30,40,50,60,70}{\draw[corAt!78,thick] (\x,0)--(\x,45);}
\foreach \x/\y in {10/0,30/0,50/0,70/0,70/15,70/30,70/45,50/45,30/45,10/45,10/30,10/15}{\fill[corNorma] (\x,\y) circle (0.18);}
\draw[corNorma,very thick,dashed] (1.0,31)--(1.0,29)--(10,29)--(10,0)--(70,0);
\node[arqSub,text=corNorma] at (30,1.6){malha de aterramento / equipotencialização};
\node[arqSub,text=corAt!80!black] at (39,43.2){captação conceitual — definir pela análise de risco};
\end{tikzpicture}
\vspace{0.3mm}\scriptsize SPDA, estruturas, BEP, blindagens, aterramento da subestação e DPS são coordenados como um único sistema equipotencial.
\end{frame}

\begin{frame}{Industrial II — normas da interface com a concessionária}
\small
\begin{columns}[T,onlytextwidth]
\column{0.49\textwidth}\begin{caixaNorma}[title=Equatorial]\begin{itemize}\item NT.00002.EQTL vigente para fornecimento em média tensão;\item padrões de medição/proteção e materiais complementares vigentes;\item análise do ponto de conexão e dados da rede antes do executivo.\end{itemize}\end{caixaNorma}
\column{0.49\textwidth}\begin{caixaProjeto}[title=Instalação do cliente]\begin{itemize}\item ABNT NBR 14039 na MT;\item ABNT NBR 5410 na BT;\item série ABNT NBR 5419 para SPDA;\item estudos de curto, seletividade, aterramento e segurança.\end{itemize}\end{caixaProjeto}
\end{columns}
\end{frame}
